\documentclass{article}
\usepackage{/globals/code}
\usepackage{/globals/anki}
\ankishow
\ankiprefix{helloworld}
\ankitopic{HelloWorld}
\begin{document}
\section{Hello World!} 
Das Ziel des Programmierens ist es einem Computer gezielt für diesen verständliche Anweisungen zu geben, so dass diese eine bestimmte Aufgabe erledigen kann. Die Sammlung dieser Anweisungen kann ein Algorithmus genannt werden.

In diesem Modul geschieht dies in Java.

\subsection{Aufbau eines Java Programms}
Der Quellcode eines Java Programms wird in einer \texttt{.java} geschrieben. Im Sinne dieses Moduls ist diese Datei das ganze Programm, und kann auch den Programmnamen als Dateinamen haben. In dieser Datei wird eine Klasse deklariert, welche auch im Sinne dieses Modules den Programmnamen als Klassennamen nutzen kann.
\begin{slsourcecode}{java}
public class HelloWorld
\end{slsourcecode}
Mit geschweiften Klammern wird ein Block angefangen. Dort drin muss eine Hauptmethode definiert werden; die Methode mit der das Programm gestartet wird.
\begin{slsourcecode}{java}
public static void main(String[] args)
\end{slsourcecode}
Diese startet auch einen Block, in welchem Anweisungen gegeben werden könnn.

\ankinote{class}{Klasse}{public class HelloWorld}
\ankinote{main}{Main}{public static void main(String[] args)}

\subsection{Konsolenausgaben}
Mit \texttt{System.out.println("...")} kann der Text der in die Klammern in Anführungszeichen geschrieben wird in die Konsole ausgegeben werden.
\begin{slsourcecode}{java}
System.out.println("Hello World");
\end{slsourcecode}
Das \texttt{ln} in \texttt{println} steht hier für line; es für automatisch ein Zeilenumbruch hinzugefügt. Wird nur \texttt{print} genutzt, passiert dies nicht.

\ankinote{print}{Print}{System.out.println("Hello World");}

\subsection{Ein komplettes Hello World}
\begin{sourcecode}{java}
public class HelloWorld {
    public static void main(String[] args) {
        System.out.println("Hello World!");
    }
}
\end{sourcecode}

\subsection{Kommentare}
Der Teil einer Zeile hinter zwei Schrägstrichen ist ein Kommentar. Ein mehrzeiliger Kommentar kann mit einem Schrägstrich und einem Sternchen angefangen werden und andersherum beendet werden. Dokumentationskommentare sind mehrzeilige Kommenrate mit einem zweitern Sternchen am anfang.
\begin{sourcecode}{java}
System.out.println("Hello World!"); // Test
/* Langer
   Kommentar*/
/** Dokumentations
   Kommentar*/
\end{sourcecode}
\end{document}